\section*{अध्याय \ref{ch:g10:soil-formation} --- मृदा, धीरे-धीरे बनने वाला संसाधन}

\begin{solution}{exo:g10:soil-formation:1}
\omterm{def:g10:soil-formation:chemical-weathering}{रासायनिक अपक्षय} जल और उसके द्वारा ले जाई जाने वाली गैसों (कार्बन डाइऑक्साइड,
ऑक्सीजन) के साथ अभिक्रियाओं द्वारा किसी चट्टान के खनिजों का रूपांतरण है: कुछ
खनिज घुल जाते हैं, कुछ सतह पर स्थायी नए खनिज बन जाते हैं।
\omterm{def:g10:soil-formation:hydrolysis}{जल-अपघटन} जल और उसके हाइड्रोजन आयनों द्वारा किसी सिलिकेट का \omterm{def:g10:soil-formation:chemical-weathering}{रासायनिक अपक्षय} है:
धातु आयन विलयन में चले जाते हैं और खनिज का बाक़ी भाग
एक \omterm{def:g10:soil-formation:clay}{मृत्तिका खनिज} में पुनर्गठित हो जाता है।
\end{solution}

\begin{solution}{exo:g10:soil-formation:2}
\omterm{def:g10:soil-formation:weathering-profile}{सैप्रोलाइट} अपने स्थान पर ही अपघटित चट्टान है: नरम और मिट्टी जैसी होते हुए भी वह अपनी
\omterm{def:g10:soil-formation:weathering-profile}{जनक चट्टान} के कण, पट्टियाँ और संरचना बनाए रखती है और नीचे की ओर, बिना किसी विराम के,
ताज़ी चट्टान में बदल जाती है। नदी की रेत में बनी मृदा कहीं और से लाए गए
पदार्थ पर टिकी होती है: उसके कण गोल और छँटे हुए होते हैं, और उसमें अपने नीचे की
चट्टान की कोई संरचना नहीं दिखती।
\end{solution}

\begin{solution}{exo:g10:soil-formation:3}
फ़ेल्डस्पार मृत्तिका (केओलिनाइट) बन जाते हैं और पोटैशियम, सोडियम या
कैल्शियम तथा सिलिका को जल में मुक्त करते हैं; काला अभ्रक मृत्तिका और लोहे के
ऑक्साइड बन जाता है; क्वार्ट्ज़ अपरिवर्तित, रेत के रूप में, बचा रहता है।
\end{solution}

\begin{solution}{exo:g10:soil-formation:4}
गहरे रंग के खनिजों का लोहा, \ce{Fe^2+}, जल द्वारा लाई गई ऑक्सीजन से
अविलेय लोहा(III) ऑक्साइडों में ऑक्सीकृत हो जाता है: पीला-भूरा गोएथाइट और लाल
हेमेटाइट, जो कणों पर परत चढ़ाकर पूरे पदार्थ को रंग देते हैं।
\end{solution}

\begin{solution}{exo:g10:soil-formation:5}
\qty{0.017}{mm} प्रति वर्ष और \qty{1.537}{mm} प्रति वर्ष: $1.537/0.017 \approx 90$
गुना बड़ी।
\end{solution}

\begin{solution}{exo:g10:soil-formation:6}
कैल्शियम: 1 और 1। कार्बन: $1 + 1 = 2$ और 2। ऑक्सीजन: $3 + 2 + 1 = 6$ और
$2 \times 3 = 6$। हाइड्रोजन: 2 और 2। आवेश: 0, और $+2 - 2 = 0$।
\end{solution}

\begin{solution}{exo:g10:soil-formation:7}
कैल्साइट के विलयन से केवल घुले हुए आयन, \ce{Ca^2+} और
\ce{HCO3-}, बनते हैं: सब कुछ बह जाता है और शुद्ध चूना पत्थर कोई ठोस
अवशेष नहीं छोड़ता। फ़ेल्डस्पार के \omterm{def:g10:soil-formation:hydrolysis}{जल-अपघटन} से एक ठोस, केओलिनाइट, बनता है, जो बना रहता है;
और ग्रेनाइट का क्वार्ट्ज़ लगभग अभिक्रिया ही नहीं करता। इसलिए ग्रेनाइट
मृत्तिका और रेत छोड़ता है।
\end{solution}

\begin{solution}{exo:g10:soil-formation:8}
मृत्तिका की पपड़ियों और ह्यूमस पर ऋणात्मक आवेश होते हैं जो धनात्मक
पोटैशियम आयनों \ce{K+} को अपनी सतहों पर थामे रखते हैं, ताकि बहकर निकलता वर्षा-जल
उन्हें बहा न ले जाए। एक जड़ हाइड्रोजन आयन छोड़ती है, जो पपड़ियों पर पोटैशियम का
स्थान ले लेते हैं; मृदा-जल में मुक्त हुआ पोटैशियम
जड़ द्वारा ग्रहण कर लिया जाता है।
\end{solution}

\begin{solution}{exo:g10:soil-formation:9}
\omterm{def:g10:soil-formation:chemical-weathering}{रासायनिक अपक्षय} को जल चाहिए, जो आर्द्र जलवायु में प्रचुर और शुष्क जलवायु में दुर्लभ होता है
(और ठंडी जलवायु में जमा हुआ); और उसकी अभिक्रियाएँ गरमी में अधिक तेज़ चलती
हैं। केवल ऊष्मा और प्रचुर जल के साथ, लंबे समय तक, ही अभिक्रियाएँ
इतनी आगे जाती हैं कि लगभग सब कुछ विलेय हट जाए और लोहे तथा
ऐलुमिनियम के ऑक्साइड बच जाएँ।
\end{solution}

\begin{solution}{exo:g10:soil-formation:10}
\qty{1}{cm} के लिए $10/0.017 \approx 590$ वर्ष; \qty{30}{cm} के लिए $300/0.017 \approx 17\,600$
वर्ष; \qty{1}{m} के लिए $1000/0.017 \approx 59\,000$ वर्ष।
\end{solution}

\begin{solution}{exo:g10:soil-formation:11}
आयतन: $10\,000 \times 0.0015 = \qty{15}{m^3}$; द्रव्यमान: $15 \times 1200 =
\qty{18000}{kg}$, प्रति हेक्टेयर प्रति वर्ष \num{18}~टन।
\end{solution}

\begin{solution}{exo:g10:soil-formation:12}
$46\,000/0.017 \approx 2.7$ मिलियन वर्ष। इतनी मोटी \omterm{def:g10:soil-formation:weathering-profile}{सैप्रोलाइट} के लिए ऐसा
भूदृश्य चाहिए जो कई मिलियन वर्षों तक स्थिर रहा हो, और जिसका अपक्षयित आवरण
उतारा न गया हो।
\end{solution}

\begin{solution}{exo:g10:soil-formation:13}
गोएथाइट: \ce{4 Fe^2+ + O2 + 6 H2O -> 4 FeO(OH) + 8 H+}: लोहा 4 और 4,
ऑक्सीजन $2 + 6 = 8$ और $4 \times 2 = 8$, हाइड्रोजन 12 और $4 + 8 = 12$, आवेश
$+8$ और $+8$। हेमेटाइट: \ce{4 Fe^2+ + O2 + 4 H2O -> 2 Fe2O3 + 8 H+}: लोहा 4
और 4, ऑक्सीजन $2 + 4 = 6$ और 6, हाइड्रोजन 8 और 8, आवेश $+8$ और $+8$।
\end{solution}

\begin{solution}{exo:g10:soil-formation:14}
शुद्ध हानि $0.082 - 0.017 = \qty{0.065}{mm}$ प्रति वर्ष; $300/0.065 \approx 4600$
वर्ष। पारंपरिक जुताई में $300/1.52 \approx 200$ वर्ष: बीस गुना से भी
अधिक कम।
\end{solution}

\begin{solution}{exo:g10:soil-formation:15}
एक टन में \qty{300}{kg} ऑर्थोक्लेज़ है। हर किलोग्राम
\qty{464}{g} केओलिनाइट देता है: $300 \times 0.464 \approx \qty{139}{kg}$। पोटैशियम:
$300 \times 0.1405 \approx \qty{42}{kg}$।
\end{solution}

\begin{solution}{pb:g10:soil-formation:1}
\textbf{1.} $4 \times 0.3048 \approx \qty{1.22}{m}$।

\textbf{2.} $0.5 \times 30.48 \approx \qty{15}{cm}$।

\textbf{3.} ह्यूमस, वह गहरे रंग का कार्बनिक पदार्थ जिसे अपघटक मृत
पौधों से बनाते हैं।

\textbf{4.} लोहे के ऑक्साइड, गोएथाइट और हेमेटाइट, जो जनक पदार्थ के खनिजों के
लोहा(II) के ऑक्सीकरण से बने।

\textbf{5.} नहीं। नीचे प्रोफ़ाइल अपक्षयित चट्टान, यानी
\omterm{def:g10:soil-formation:weathering-profile}{सैप्रोलाइट} (या कोई अन्य अपक्षयित पदार्थ), में और फिर ताज़ी \omterm{def:g10:soil-formation:weathering-profile}{जनक
चट्टान} में बदल जाएगा।

\textbf{6.} \ce{2 KAlSi3O8 + 2 H+ + 9 H2O -> Al2Si2O5(OH)4 + 2 K+ + 4 H4SiO4}।

\textbf{7.} $1000/278.33 \approx \qty{3.59}{mol}$।

\textbf{8.} $3.59/2 \approx \qty{1.80}{mol}$ केओलिनाइट, $1.80 \times 258.16
\approx \qty{464}{g}$।

\textbf{9.} पोटैशियम आयन और घुली हुई सिलिका: जल उन्हें नदियों तक ले जाता है,
और नदियाँ समुद्र तक।

\textbf{10.} वह लगभग अभिक्रिया नहीं करता: क्वार्ट्ज़ के कण रेत के रूप में बचे रहते हैं।

\textbf{11.} $1219/0.017 \approx 72\,000$ वर्ष।

\textbf{12.} $1.537 - 0.017 = \qty{1.52}{mm}$ प्रति वर्ष।

\textbf{13.} $1219/1.52 \approx 800$ वर्ष।

\textbf{14.} $72\,000/800 = 90$: प्रोफ़ाइल जितनी तेज़ी से बना उससे नब्बे गुना तेज़ी से
खो जाता है।

\textbf{15.} नहीं: $0.013 < 0.017$, अपरदन उत्पादन से थोड़ा धीमा है
और मृदा टिकी रहती है।

\textbf{16.} $T = S/(E - P)$।

\textbf{17.} $300/1.52 \approx 200$ वर्ष।

\textbf{18.} $300/(0.082 - 0.017) = 300/0.065 \approx 4600$ वर्ष।

\textbf{19.} कुछ शताब्दियों में खोई मृदा को फिर से बनने में दसियों हज़ार वर्ष
लगते हैं: मानव जीवन और समाजों के पैमाने पर जो खो गया उसकी
पूर्ति नहीं होती।

\textbf{20.} $1000/1.52 \approx 660$ वर्ष: \textbf{पारंपरिक जुताई में एक मीटर मृदा
लगभग 660 वर्षों में, यानी कुछ शताब्दियों में, खो जाती है}।
\end{solution}
